We find imperfections in the normal approximation but little sensitivity to truncation. Among 10,476 matched records, the median absolute mean--median gap is 0.11 estimated standard deviations and the median fitted probability outside 5--10 is about 0.04\%. The extreme tail fits less well: 23\% of available maxima lie outside their pointwise 95\% predictive intervals, and observed maxima exceed their expected values by 0.17 grade points on average. These discrepancies suggest that a normal distribution does not fully describe the upper tail; they are not formal evidence of a particular alternative distribution.

When we truncated the normal distributions to $[5,10]$, keeping the same underlying parameters and graduate weights, selection shares changed by a graduate-weighted average absolute difference of only 0.04 percentage points. The largest change in a field's average was 0.13 percentage points, and no field changed its overall rank. A pooled top-20\% threshold concerns the upper fifth rather than the single largest grade, explaining why imperfect maxima can coexist with stable selection estimates. We therefore regard Figure~\ref{fig:selection} as robust to this truncation check, without claiming robustness to every possible distribution or establishing normality from summaries.
